% ============================================================================
% MODELO METALINGUÍSTICO — Relatório Técnico de Extensão como Prática Profissional
% Curso Técnico Integrado em Informática (PPC 2012) — IFRN Campus Ceará-Mirim
% ============================================================================
%
% O QUE É ESTE ARQUIVO
% Cada seção deste modelo traz, dentro de um quadro em itálico, a explicação
% do que essa seção deve conter (aparece no PDF gerado — não é comentário
% oculto). Ao escrever o relatório de verdade, apague o quadro de instruções
% e escreva o texto no lugar dele.
% Para ver como fica um relatório já preenchido, com esse mesmo modelo, veja
% o arquivo irmão "relatorio-exemplo-ficticio.tex".
%
% BASE NORMATIVA (resumida; ver README.md desta pasta e de ../ para os
% artigos citados)
%  - OD 2024, Art. 291, I, "a": projeto de extensão registrado por relatório
%    técnico (mesmo caminho do projeto de pesquisa/integrador — não há
%    modalidade "artigo científico" para extensão).
%  - OD 2024, Art. 296: ABNT para redação de trabalhos técnicos e científicos.
%  - OD 2024, Art. 305: relatório produzido com base no projeto desenvolvido,
%    no período de realização do projeto.
%  - OD 2024, Art. 304: avaliação por banca (orientador(a) + professores das
%    disciplinas vinculadas), nota de 0 a 100, aprovação com mínimo de 60.
%  - OD 2024, Art. 300, §2º: prazo de 30 dias após concluir o projeto, no
%    técnico integrado, para entregar o documento ao(à) orientador(a).
%  - PPC 2012, seção 5.2.1: unidade entre teoria e prática, aplicação dos
%    conhecimentos do curso, intervenção no mundo do trabalho e na realidade
%    social — para extensão, isso significa mostrar a que necessidade
%    concreta da comunidade/instituição parceira o projeto respondeu.
%  - Guia Prático para Relatório Técnico ou Científico (Biblioteca do IFRN,
%    2011, NBR 10719:2011): estrutura e formatação usadas neste modelo.
%
% Compila com pdflatex (rode duas vezes por causa de referências internas).
% ============================================================================

\documentclass[article,12pt,a4paper]{abntex2}

\usepackage[T1]{fontenc}
\usepackage[utf8]{inputenc}
\usepackage[brazil]{babel}
\usepackage{times}
\usepackage{indentfirst}
\usepackage{graphicx}
\usepackage{microtype}

% Quadro de instruções: tudo que estiver dentro deste ambiente é explicação
% de apoio, não texto do relatório. Apague o ambiente inteiro (e o que está
% dentro) ao escrever o conteúdo de verdade.
\newenvironment{instrucoes}{%
  \begin{quotation}\footnotesize\itshape
  \noindent\textbf{\upshape Instruções desta seção — apague antes de entregar:}\par\smallskip
}{%
  \end{quotation}
}

% ---- Dados do trabalho (edite aqui) ----------------------------------------
\renewcommand{\maketitlehookb}{}%% sem título em língua estrangeira (sem abstract)
\titulo{[TÍTULO DO RELATÓRIO — nome da ação de extensão + instituição/comunidade atendida]}
\autor{[Nome do(a) estudante] \and [Nome do(a) coautor(a), se houver] \and [Nome do(a) orientador(a)]\footnote{Professor(a) orientador(a). [Titulação e lotação].}}
\local{Ceará-Mirim -- RN}
\data{[ano]}
\instituicao{Instituto Federal de Educação, Ciência e Tecnologia do Rio Grande do Norte -- IFRN \\ Campus Ceará-Mirim}

\begin{document}
\maketitle

% =====================================================================
\begin{resumo}
\begin{instrucoes}
Regras do resumo (NBR 6028, conforme o Guia do IFRN) — mesmas de qualquer relatório técnico:
\begin{itemize}
  \item Elemento obrigatório. Sequência de frases concisas, nunca uma lista de tópicos.
  \item Parágrafo único, de 150 a 500 palavras, espaçamento 1,5.
  \item Ordem fixa: objetivo $\to$ metodologia $\to$ resultados $\to$ conclusão. Para extensão, "objetivo" é a necessidade atendida e "resultados" é o benefício gerado ao público.
  \item Verbo na voz ativa e na 3\textsuperscript{a} pessoa do singular.
  \item Evitar símbolos, fórmulas, abreviações e citações.
  \item Ao final, "Palavras-chave" em negrito, seguidas de 3 a 5 termos separados por ponto e finalizados por ponto.
\end{itemize}
\end{instrucoes}

[Escreva aqui o resumo, seguindo as regras acima.]

\vspace{\onelineskip}
\noindent
\textbf{Palavras-chave}: [Termo 1]. [Termo 2]. [Termo 3]. [Termo 4, opcional].
\end{resumo}

\textual

% =====================================================================
\section{Introdução}

\begin{instrucoes}
Esta seção precisa cumprir cinco pontos, adaptados ao caráter extensionista do projeto:
\begin{enumerate}
  \item \textbf{Contextualização}: qual é a instituição, escola, comunidade ou público que o projeto atendeu, e qual necessidade concreta dele motivou a ação — é o que caracteriza a "intervenção no mundo do trabalho e na realidade social" exigida pelo PPC (seção 5.2.1) para valer como prática profissional.
  \item \textbf{Justificativa}: por que essa necessidade merecia ser atendida por um projeto de extensão do curso, e não resolvida de outra forma; qual a relação com a área técnica do curso (Informática).
  \item \textbf{Objetivo geral} (uma frase: o que o projeto se propôs a entregar ao público atendido) e, se fizer sentido, \textbf{objetivos específicos}.
  \item \textbf{Estrutura do relatório}: uma frase final dizendo o que vem em cada seção seguinte.
  \item Conexão com o \textbf{plano de atividades} deferido pelo(a) orientador(a) (Res. 25/2019, Art. 4º, IV) — é o que a banca vai comparar na avaliação (OD 2024, Art. 304).
\end{enumerate}
\end{instrucoes}

[Escreva aqui a introdução, cobrindo os cinco pontos acima.]

% =====================================================================
\section{Fundamentação}

\begin{instrucoes}
Mais breve do que num relatório de pesquisa. Deve conter:
\begin{itemize}
  \item Os conceitos técnicos necessários para entender as atividades realizadas (ex.: os fundamentos de manutenção de hardware, redes, ou o que for pertinente ao serviço prestado), explicados na profundidade adequada a um trabalho de curso técnico.
  \item Uma breve nota sobre o papel da extensão no IFRN: articulação entre ensino, pesquisa e extensão, e via de mão dupla entre a instituição e a comunidade (não é obrigatório citar autor para isso, mas se citar algum documento institucional, referenciar corretamente).
  \item Se houver literatura técnica ou documentação de apoio consultada, citar conforme NBR 10520 (autor-data) e listar depois em Referências.
  \item Encerrar relacionando os conceitos com os conteúdos das disciplinas do núcleo tecnológico do curso.
\end{itemize}
\end{instrucoes}

[Escreva aqui a fundamentação, com citações no formato autor-data quando houver, por
exemplo: "Segundo Silva (2020), ..." ou "(SILVA, 2020, p. 12)".]

% =====================================================================
\section{Metodologia e Descrição das Atividades}

\begin{instrucoes}
Deve responder, nesta ordem:
\begin{itemize}
  \item \textbf{Planejamento}: como o projeto foi organizado (cronograma, etapas), e como o público-alvo/instituição parceira foi definido e contatado.
  \item \textbf{Público atendido}: quem foi beneficiado (número de pessoas, turmas, equipamentos, instituições), e como isso foi levantado.
  \item \textbf{Procedimento}: passo a passo do que foi feito em campo — atendimentos realizados, oficinas ministradas, equipamentos consertados/configurados, materiais produzidos, conforme o caso.
  \item \textbf{Materiais e ferramentas} usados, com detalhe suficiente para outra pessoa reproduzir a ação.
  \item \textbf{Local e período} de realização, compatíveis com o plano de atividades deferido pelo(a) orientador(a) e com o prazo de entrega (30 dias após concluir o projeto, no técnico integrado — OD 2024, Art. 300, §2º).
\end{itemize}
\end{instrucoes}

[Escreva aqui a metodologia e a descrição das atividades, respondendo: o quê, para quem, com quê, como e quando.]

% =====================================================================
\section{Resultados e Impacto}

\begin{instrucoes}
Deve conter:
\begin{itemize}
  \item Números concretos do que foi entregue (quantidade de equipamentos atendidos, pessoas capacitadas, chamados resolvidos, oficinas realizadas etc.), com tabela se ajudar a organizar os dados.
  \item Avaliação do impacto: o que mudou para o público atendido (ex.: retomada de uso dos equipamentos, redução de tempo de espera por suporte, depoimento da instituição parceira, quando houver).
  \item Resposta explícita ao objetivo e, se houver, a cada objetivo específico definido na introdução.
  \item Diferente de um relatório de pesquisa, aqui o foco não é "o que descobrimos", é "que diferença fizemos" — mantenha esse ângulo.
\end{itemize}
\end{instrucoes}

[Escreva aqui os resultados e o impacto. Exemplo de tabela abaixo — apague se não for usar.]

\begin{table}[!ht]
  \centering
  \caption{[Legenda da tabela, descrevendo o que ela mostra]}
  \begin{tabular}{lcc}
    \hline
    [Categoria de atendimento] & [Quantidade] & [Observação] \\
    \hline
    [Item A] & [valor] & [valor] \\{}
    [Item B] & [valor] & [valor] \\
    \hline
  \end{tabular}
  \\[0.3em]
  {\footnotesize Fonte: [registros do projeto / elaborado pelo(s) autor(es), ano].}
\end{table}

% =====================================================================
\section{Considerações Finais}

\begin{instrucoes}
Deve conter:
\begin{itemize}
  \item Retomada do objetivo: o projeto atendeu o que se propôs? Como?
  \item Síntese dos principais resultados, sem repetir os números já mostrados em detalhe na seção anterior.
  \item Contribuição do projeto: o que ele agregou à comunidade/instituição parceira, e o que agregou à formação técnica do(a) estudante — este é o ponto em que se demonstra a "articulação entre teoria e prática" que caracteriza a prática profissional (PPC 2012, seção 5.2.1).
  \item Limitações do trabalho (o que não foi possível atender, e por quê).
  \item Continuidade: se a ação pode/deve continuar, e como (útil se a instituição parceira for manter vínculo com o curso).
  \item Não introduzir dado novo aqui: esta seção só retoma o que já foi mostrado.
\end{itemize}
\end{instrucoes}

[Escreva aqui as considerações finais.]

\postextual

% =====================================================================
\section*{Referências}

\begin{instrucoes}
Regras (NBR 6023):
\begin{itemize}
  \item Ordem alfabética pelo sobrenome do primeiro autor.
  \item Um parágrafo por obra, sem numeração, com recuo da segunda linha em diante.
  \item Todo item citado no texto precisa estar aqui, e só o que foi citado deve estar aqui. Se o relatório não citar nenhuma fonte externa (comum em relatórios de extensão mais operacionais), esta seção pode ficar só com a documentação técnica de apoio usada (manuais, datasheets, documentação oficial de software).
\end{itemize}
Exemplos de formato por tipo de fonte (adapte ao que você usou):
\end{instrucoes}

\newcommand{\referencia}[1]{\par\noindent\hangindent=1cm #1\par\medskip}

\referencia{SOBRENOME, Nome. \textit{Título da obra}: subtítulo se houver.
Edição. Local de publicação: Editora, ano.}

\referencia{NOME DA INSTITUIÇÃO/FABRICANTE. \textit{Título do manual ou
documentação técnica}. Local: Instituição, ano. Disponível em:
\texttt{https://endereco-completo}. Acesso em: dia mês ano.}


\end{document}
